\documentclass[11pt]{article}
\usepackage[margin=1in]{geometry}
\usepackage{amsmath,amssymb,amsthm}
\usepackage{xurl}
\usepackage{hyperref}
\sloppy
\let\oldus\_ \renewcommand{\_}{\oldus\allowbreak}
\newtheorem{theorem}{Theorem}
\newtheorem{lemma}[theorem]{Lemma}
\theoremstyle{definition}
\newtheorem{definition}[theorem]{Definition}
\theoremstyle{remark}
\newtheorem{remark}[theorem]{Remark}

\title{Disproof of conjecture 00000003524:\\ $\omega_1$ has two disjoint stationary subsets}
\author{Jackmeson1}
\date{2026-10-05}

\begin{document}
\maketitle

\section{The conjecture and the clause refuted}

\textbf{Statement (English, verbatim).} Definition: Subsets of $\omega_1$: the combinatorics of
uncountable ordinals. Conjecture: Club sets and stationary sets: the intersection of stationary
sets is stationary, with the reflection of stationarity bounded below closure. (stationary
intersection reflection lower bound closure)

\textbf{Statement (Chinese).} The Chinese text says, word for word: ``Definition: subsets of
$\omega_1$: the combinatorics of uncountable ordinals. Conjecture: club sets and stationary sets:
the intersection of stationary sets is stationary (\emph{pingwen ji de jiaoji wei pingwen}) and
(\emph{qie}) the reflection of stationarity is the lower-bound closure.''

\textbf{Reading.} Both languages literally assert
\[
  \text{(A)}\qquad S, T \subseteq \omega_1 \text{ stationary} \implies S\cap T \text{ stationary},
\]
joined to a second clause (B), ``the reflection of stationarity bounded below closure''. The
Chinese text joins (A) and (B) with ``and'' (\emph{qie}), so the conjecture is the conjunction
(A) $\wedge$ (B). Clause (B) does not state a definite proposition: stationary reflection is a
standard notion, but ``bounded below closure'' (``is the lower-bound closure'') has no standard
meaning. We do not interpret (B). We refute (A), and we prove in Lean that (A) $\wedge$ $Q$ is
false for \emph{every} proposition $Q$, so the conjunction fails however (B) is read.

We refute (A) in three forms: for two sets, for every nonempty family of stationary sets, and
for the family of all stationary sets. We do \emph{not} refute the different statement ``a
stationary set intersected with a club set is stationary''. That statement is true, and we
prove it in Lean as a remark (Remark~\ref{rem:club}). The text says ``the intersection of
stationary sets'', not ``of a stationary set and a club set''.

\section{Definitions (as formalized)}

We use Mathlib's own definitions from
\texttt{Mathlib/SetTheory/Cardinal/Cofinality/Club.lean}, unchanged.

\begin{definition}
$\omega_1$ is Mathlib's \texttt{Ordinal.omega 1}, the first uncountable ordinal. Subsets of
$\omega_1$ are subsets of the type \texttt{($\omega_1$).ToType}, the well-order of order type
$\omega_1$ whose elements correspond to the countable ordinals.
\end{definition}

\begin{definition}[club]
In a linear order $\alpha$, a set $C$ is a \emph{club} (\texttt{IsClub C}) if
(i) it is closed: whenever $d\subseteq C$ is nonempty and $a$ is a least upper bound of $d$,
then $a\in C$; and (ii) it is unbounded (cofinal): for every $x$ there is $y\in C$ with
$x\le y$. Mathlib states (i) as \texttt{DirSupClosed C}, which also asks that $d$ be directed.
Every subset of a linear order is directed, so this is the same condition. The Lean theorem
\texttt{isClub\_iff\_closed\_unbounded} proves the equivalence of Mathlib's \texttt{IsClub} with
(i) $\wedge$ (ii). In $\omega_1$ a set with a least upper bound is bounded, so (i) says that $C$
contains the supremum of each of its nonempty bounded subsets. This is the usual notion of a
closed subset of $\omega_1$.
\end{definition}

\begin{definition}[stationary]
$S\subseteq\alpha$ is \emph{stationary} (\texttt{IsStationary S}) if $S\cap C\ne\emptyset$ for
every club $C$. The Lean theorem \texttt{isStationary\_iff\_meets\_every\_club} states this
unfolding. In particular $\emptyset$ is not stationary (Mathlib's
\texttt{not\_isStationary\_empty}), because $\alpha$ itself is a club.
\end{definition}

\section{Proof}

Let $\alpha$ be an uncountable linear order with a well-founded $<$, in which every initial
segment $\{x : x < a\}$ is countable and whose cofinality \texttt{Order.cof} $\alpha$ is not
$\aleph_0$. All three hypotheses hold for $\alpha = \omega_1$:
$|\omega_1| = \aleph_1 > \aleph_0$; each initial segment has cardinality $<\aleph_1$, hence is
countable (Mathlib's \texttt{mk\_Iio\_lt}); and $\operatorname{cof}\omega_1 = \aleph_1$ (Mathlib's
\texttt{Ordinal.cof\_toType}, \texttt{Cardinal.cof\_omega\_one}).

\begin{lemma}[\texttt{isStationary\_Ioi}]\label{lem:tail}
For every $b\in\alpha$ the final segment $(b,\to) = \{a : b < a\}$ is stationary.
\end{lemma}
\begin{proof}
If its complement $\{x : x\le b\}$ were cofinal, every $x$ would satisfy $x\le b$, so
$\alpha = \{x : x<b\}\cup\{b\}$ would be countable. Hence the complement is not cofinal, and a
set whose complement is not cofinal is stationary (Mathlib's
\texttt{IsStationary.of\_not\_isCofinal\_compl}).
\end{proof}

\begin{theorem}[Ulam matrix; \texttt{exists\_disjoint\_stationary}]\label{thm:ulam}
$\alpha$ contains stationary sets $S$, $T$ with $S\cap T=\emptyset$.
\end{theorem}
\begin{proof}
For each $a\in\alpha$ the initial segment below $a$ is countable, so we can choose
$f_a:\alpha\to\mathbb{N}$ injective on $\{x : x<a\}$. For $n\in\mathbb{N}$ and $b\in\alpha$ put
\[
  A(n,b) = \{a\in\alpha : b<a \ \text{and}\ f_a(b)=n\}.
\]
For fixed $b$, $\bigcup_{n\in\mathbb{N}} A(n,b) = (b,\to)$, which is stationary by
Lemma~\ref{lem:tail}. Since $\operatorname{cof}\alpha\neq\aleph_0$, a countable union of sets
is stationary if and only if one of the sets is (Mathlib's
\texttt{isStationary\_iUnion\_iff\_of\_countable}). So there is $N(b)\in\mathbb{N}$ with
$A(N(b),b)$ stationary. The map $N:\alpha\to\mathbb{N}$ is not injective, because $\alpha$ is
uncountable. So there are $b\ne b'$ with $N(b)=N(b')=n$. Put $S=A(n,b)$ and $T=A(n,b')$. If
$a\in S\cap T$, then $b,b'<a$ and $f_a(b)=n=f_a(b')$. Since $f_a$ is injective on
$\{x : x<a\}$, this gives $b=b'$, a contradiction. Hence $S\cap T=\emptyset$.
\end{proof}

\begin{theorem}[\texttt{stationaryInterClaim\_false}, \texttt{sInter\_claim\_false},
\texttt{not\_isStationary\_sInter\_all}, \texttt{conjecture3524\_false}]
In $\omega_1$:
(1) there are stationary $S$, $T$ with $S\cap T$ not stationary, so (A) is false;
(2) it is false that every nonempty family of stationary sets has a stationary intersection;
(3) the intersection of all stationary subsets of $\omega_1$ is not stationary;
(4) for every proposition $Q$, $\neg(\text{(A)}\wedge Q)$.
\end{theorem}
\begin{proof}
Take $S$, $T$ from Theorem~\ref{thm:ulam} with $\alpha=\omega_1$. Then $S\cap T=\emptyset$, and
$\emptyset$ is not stationary. This gives (1). For (2) apply the hypothesis to the family
$\{S,T\}$, and for (3) note that the intersection of all stationary sets lies inside $S\cap T$,
so it is empty. (4) follows from (1).
\end{proof}

\begin{remark}[\texttt{isStationary\_inter\_isClub}; not part of the refutation]\label{rem:club}
If $S\subseteq\omega_1$ is stationary and $C\subseteq\omega_1$ is a club, then $S\cap C$ is
stationary: for any club $D$, $C\cap D$ is a club (Mathlib's \texttt{IsClub.inter}), so
$(S\cap C)\cap D = S\cap(C\cap D)\ne\emptyset$. This is a different statement from (A).
\end{remark}

\section{Lean formalization and axiom audit}

The file \texttt{lean/Conjecture3524/Basic.lean} (175 lines, \texttt{import Mathlib} only,
namespace \texttt{C3524}, Lean v4.33.1, Mathlib v4.33.1) contains:
\begin{itemize}
\item \texttt{isClub\_iff\_closed\_unbounded}, \texttt{isStationary\_iff\_meets\_every\_club}:
  the unfoldings of Mathlib's definitions given in Section 2.
\item \texttt{isStationary\_Ioi}, \texttt{exists\_disjoint\_stationary}: Lemma~\ref{lem:tail} and
  Theorem~\ref{thm:ulam} for a general linear order $\alpha$.
\item \texttt{countable\_Iio\_omega\_one}, \texttt{cof\_omega\_one\_toType},
  \texttt{uncountable\_omega\_one\_toType}: the hypotheses for $\omega_1$.
\item \texttt{exists\_disjoint\_stationary\_omega\_one}:
  \texttt{$\exists$ S T : Set ($\omega_1$).ToType, IsStationary S $\wedge$ IsStationary T
  $\wedge$ Disjoint S T}.
\item \texttt{StationaryInterClaim} (the literal clause (A)):
  \texttt{$\forall$ S T : Set ($\omega_1$).ToType, IsStationary S $\to$ IsStationary T $\to$
  IsStationary (S $\cap$ T)}, and \texttt{stationaryInterClaim\_false :
  $\neg$ StationaryInterClaim}.
\item \texttt{sInter\_claim\_false}, \texttt{not\_isStationary\_sInter\_all}: forms (2) and (3).
\item \texttt{conjecture3524\_false (Q : Prop) : $\neg$ (StationaryInterClaim $\wedge$ Q)}:
  the main theorem.
\item \texttt{isStationary\_inter\_isClub}: Remark~\ref{rem:club}.
\end{itemize}

\textbf{Conventions.} The statements are universe-polymorphic: $\omega_1$ is taken in
\texttt{Ordinal.\{u\}} for an arbitrary universe \texttt{u}. Choice is used (via
\texttt{Classical.choice}) to pick the injections $f_a$ and the indices $N(b)$.

\textbf{Axiom audit.} \texttt{\#print axioms} for \texttt{conjecture3524\_false},
\texttt{stationaryInterClaim\_false}, \texttt{sInter\_claim\_false},
\texttt{not\_isStationary\_sInter\_all}, \texttt{exists\_disjoint\_stationary\_omega\_one} and
\texttt{isStationary\_inter\_isClub} reports \texttt{[propext, Classical.choice, Quot.sound]}.
The file elaborates with no errors and no warnings. It contains no \texttt{sorry}, no
\texttt{native\_decide} and no added axioms.

\section*{Sources}
\begin{itemize}
\item Wikipedia, ``Stationary set'',
  \url{https://en.wikipedia.org/wiki/Stationary_set}: ``a stationary set is a set that is not too
  small in the sense that it intersects all club sets''; for $\kappa$ of countable cofinality
  ``any two stationary subsets of $\kappa$ have stationary intersection. This is no longer the
  case if the cofinality of $\kappa$ is uncountable.'' A stationary subset of a regular $\kappa$
  ``can be partitioned into $\kappa$ many disjoint stationary sets. This result is due to
  Solovay. If $\kappa$ is a successor cardinal, this result is due to Ulam and is easily shown
  by means of what is called an Ulam matrix.''
\item Wikipedia, ``Club set'', \url{https://en.wikipedia.org/wiki/Club_set}: ``If a set is both
  closed and unbounded, then it is a club set''.
\item Mathlib, \texttt{Mathlib/SetTheory/Cardinal/Cofinality/Club.lean} (definitions
  \texttt{IsClub}, \texttt{IsStationary}; lemmas \texttt{IsClub.inter},
  \texttt{IsStationary.of\_not\_isCofinal\_compl},
  \texttt{isStationary\_iUnion\_iff\_of\_countable}, \texttt{not\_isStationary\_empty}).
\end{itemize}

\end{document}
